\documentclass[11pt,a4paper]{ctexart}
\usepackage[margin=1.55cm]{geometry}
\usepackage{amsmath,amssymb,bm}
\usepackage{booktabs,longtable,array}
\usepackage{xcolor}
\usepackage{tcolorbox}
\usepackage{hyperref}

\hypersetup{colorlinks=true, linkcolor=black, urlcolor=blue}
\setlength{\parindent}{0pt}
\setlength{\parskip}{5pt}
\renewcommand{\arraystretch}{1.45}
\newcommand{\ds}{\displaystyle}

\title{\vspace{-1.2em}第6章后半：镜像法与边角绕流讲解}
\author{}
\date{}

\begin{document}
\maketitle
\vspace{-2.2em}

\begin{tcolorbox}[colback=blue!3,colframe=blue!50!black,title=对应课件范围]
对应灰色课件《第6章 理想不可压无旋流动-叠加.pdf》的第21页后半到第24页。
\end{tcolorbox}

\begin{center}
\begin{tabular}{lll}
\toprule
页码 & 内容 & 考试价值\\
\midrule
21--22 & 平面镜像法 & 高\\
23 & 双壁面镜像法 & 中高\\
23--24 & 边角绕流 $\varphi=Ar^n\cos n\theta,\ \psi=Ar^n\sin n\theta$ & 高\\
24 & 总结：势流 = 拉普拉斯方程 + 叠加 + 验证边界 & 很高\\
\bottomrule
\end{tabular}
\end{center}

\section*{一、镜像法到底在干什么}

理想流体碰到固体壁面，边界条件不是“不滑移”，而是“不穿透”：
\[
v_n=0.
\]

二维流动中，若固体壁面是一条流线，则流体速度沿着壁面走，没有法向穿透速度。因此：
\[
\boxed{\text{固壁}=\text{流线}}
\]

镜像法的目的就是构造一个“真实奇点 + 虚拟镜像奇点”的组合，使固体壁面上
\[
\psi=\text{常数}.
\]

只要能证明这一点，就满足固壁不穿透条件。

\begin{tcolorbox}[colback=yellow!6,colframe=orange!75!black,title=考试判分点]
镜像法题不能只写“放一个镜像”。标准写法必须有三步：
\[
1.\ \text{放镜像奇点}
\]
\[
2.\ \text{写总流函数 } \psi=\psi_1+\psi_2
\]
\[
3.\ \text{代入固壁，证明 } \psi=\text{常数}
\]
\end{tcolorbox}

\section*{二、镜像法基本规则}

假设固壁是 $y=0$，真实流动区域是 $y>0$。

\begin{center}
\begin{tabular}{llll}
\toprule
真实奇点 & 镜像位置 & 镜像强度 & 原因\\
\midrule
点源 & 关于固壁对称 & 同号点源 & 法向速度抵消\\
点汇 & 关于固壁对称 & 同号点汇 & 法向速度抵消\\
点涡 & 关于固壁对称 & 反号点涡 & 让壁面成为流线\\
\bottomrule
\end{tabular}
\end{center}

记忆方式：
\[
\boxed{\text{源汇同号，点涡反号}}
\]

\section*{三、例题1：点涡在水平壁面上方}

真实点涡位于 $(0,a)$，强度为 $k$，固壁为 $y=0$，实际区域为 $y>0$。

真实点涡流函数：
\[
\psi_1=-k\ln\sqrt{x^2+(y-a)^2}.
\]

镜像点涡位于 $(0,-a)$，强度相反，因此：
\[
\psi_2=+k\ln\sqrt{x^2+(y+a)^2}.
\]

总流函数：
\[
\psi=\psi_1+\psi_2
=-k\ln\sqrt{x^2+(y-a)^2}
+k\ln\sqrt{x^2+(y+a)^2}.
\]

验证固壁 $y=0$：
\[
\psi\big|_{y=0}
=-k\ln\sqrt{x^2+a^2}
+k\ln\sqrt{x^2+a^2}
=0.
\]

所以：
\[
\boxed{y=0\text{ 是流线，满足固壁不穿透条件。}}
\]

\section*{四、例题2：点源在竖直壁面右侧}

真实区域为 $x>0$，固壁是 $x=0$。点源位于 $(a,0)$，强度为 $A$。

真实点源：
\[
\psi_1=A\tan^{-1}\frac{y}{x-a}.
\]

镜像点源在 $(-a,0)$，同号：
\[
\psi_2=A\tan^{-1}\frac{y}{x+a}.
\]

总流函数：
\[
\psi
=A\tan^{-1}\frac{y}{x-a}
+A\tan^{-1}\frac{y}{x+a}.
\]

验证固壁 $x=0$：
\[
\psi\big|_{x=0}
=A\tan^{-1}\frac{y}{-a}
+A\tan^{-1}\frac{y}{a}
=0.
\]

所以：
\[
\boxed{x=0\text{ 是流线，满足固壁不穿透条件。}}
\]

\section*{五、双壁面镜像法}

若实际区域是第一象限：
\[
x>0,\qquad y>0,
\]
则有两面固壁：
\[
x=0,\qquad y=0.
\]

如果真实点源位于 $(b,a)$，为了让两面墙都成为流线，需要补三个镜像点：
\[
(b,a),\qquad (-b,a),\qquad (b,-a),\qquad (-b,-a).
\]

点源/点汇：四个都是同号。

点涡：每关于一面墙镜像一次，符号变一次。

\begin{tcolorbox}[colback=green!4,colframe=green!45!black,title=双壁面题的标准流程]
1. 画出真实奇点。\\
2. 对第一面墙镜像。\\
3. 对第二面墙镜像。\\
4. 写总流函数。\\
5. 分别代入 $x=0$ 和 $y=0$，证明两面墙上 $\psi=\text{常数}$。
\end{tcolorbox}

\section*{六、边角绕流}

课件给出的边角绕流公式是：
\[
\boxed{\varphi=A r^n\cos n\theta}
\]
\[
\boxed{\psi=A r^n\sin n\theta}
\]

其中 $\alpha$ 是夹角，$n$ 与夹角的关系为：
\[
\boxed{n=\frac{\pi}{\alpha}}.
\]

为什么这个公式能表示边角绕流？

因为边界是两条直线：
\[
\theta=0,\qquad \theta=\frac{\pi}{n}.
\]

代入流函数：
\[
\psi=A r^n\sin n\theta.
\]

当 $\theta=0$ 时：
\[
\psi=Ar^n\sin 0=0.
\]

当 $\theta=\pi/n$ 时：
\[
\psi=Ar^n\sin \pi=0.
\]

所以两条边界都是流线，满足固壁不穿透条件。

\section*{七、边角绕流速度公式}

由势函数
\[
\varphi=A r^n\cos n\theta
\]
求速度：
\[
v_r=\frac{\partial\varphi}{\partial r}
=Anr^{n-1}\cos n\theta,
\]
\[
v_\theta=\frac{1}{r}\frac{\partial\varphi}{\partial\theta}
=-Anr^{n-1}\sin n\theta.
\]

速度大小为：
\[
\boxed{V=Anr^{n-1}}.
\]

这个公式能判断角点附近速度是否发散：

\begin{center}
\begin{tabular}{lll}
\toprule
情况 & 数学表现 & 物理含义\\
\midrule
$n>1$，即 $\alpha<180^\circ$ & $r\to0$ 时 $V\to0$ & 角点速度趋于零\\
$n=1$，即 $\alpha=180^\circ$ & $V=A$ & 均匀流\\
$n<1$，即 $\alpha>180^\circ$ & $r\to0$ 时 $V\to\infty$ & 角点速度奇异\\
\bottomrule
\end{tabular}
\end{center}

\section*{八、常见夹角速查}

\begin{center}
\begin{tabular}{ccc}
\toprule
夹角 $\alpha$ & $n=\pi/\alpha$ & 速度特征\\
\midrule
$60^\circ$ & $3$ & 角点速度趋于 $0$\\
$90^\circ$ & $2$ & 角点速度趋于 $0$\\
$120^\circ$ & $3/2$ & 角点速度趋于 $0$\\
$180^\circ$ & $1$ & 均匀流\\
$270^\circ$ & $2/3$ & 角点速度趋于无穷\\
$360^\circ$ & $1/2$ & 角点速度趋于无穷\\
\bottomrule
\end{tabular}
\end{center}

\section*{九、考试满分流程}

遇到课件第21--24页相关题，按下面流程写：

\begin{tcolorbox}[colback=red!3,colframe=red!60!black,title=统一解题流程]
\[
1.\ \text{判断题型：镜像法还是边角绕流。}
\]
\[
2.\ \text{有墙面和奇点：用镜像法。}
\]
\[
3.\ \text{有夹角区域：用 } \varphi=A r^n\cos n\theta,\quad \psi=A r^n\sin n\theta.
\]
\[
4.\ \text{验证边界：}\psi=\text{常数}.
\]
\[
5.\ \text{若要求压强：用伯努利方程。}
\]
\end{tcolorbox}

压强统一公式：
\[
p+\frac12\rho V^2=p_\infty+\frac12\rho U_\infty^2,
\]
\[
\boxed{p=p_\infty+\frac12\rho(U_\infty^2-V^2)}.
\]

\section*{十、最后记忆句}

\[
\boxed{\text{势流造解不是乱凑公式，而是让边界变成流线。}}
\]

镜像法是让直墙变成流线；边角绕流公式是让两条夹角边界变成流线。

\end{document}
